\section{Why effectivity remains a theorem, not a convention}

For a morphism of a nonsingular projective variety, dynatomic effectivity can
be formulated through graph--diagonal intersection multiplicities and proved
under nondegeneracy hypotheses \citep{Hutz2010}.  The present object is not a
direct instance of that statement.  The polynomial H\'enon automorphism of
the affine plane extends birationally, rather than morphically, to the
projective plane, and $F_n$ is a further intersection with $\Fix(J)$ followed
by a mixed-axis condition.

The exact missing statement is concrete.

\begin{quote}
For every odd $n$, the intersection
$\Fix(J)\cap H^{-(n+1)/2}\Fix(R)$ is transverse after lower minimal periods
are removed, or its local intersection multiplicities have nonnegative
M\"obius transform.
\end{quote}

A repeated root of $F_n$ is the natural finite falsifier.  Differentiating the
symmetry-line shooting equation should translate it into tangency of the two
lines, equivalently a constrained monodromy degeneracy.  The certified local
hyperbolicity may exclude that degeneracy for survivor roots, but it cannot
silently certify ambient algebraic roots.  This monodromy/tangency theorem is
the next minimal project.

Even after effectivity, a Galois-height pressure would remain separate: it
must weight all embeddings and prove a uniform asymptotic, not simply count
them.
